% =============================================================================
% Semana 04 — Infográficos de Função + Estatística Descritiva II: Variabilidade
% Aula de 3 horas (19h–22h) | 04/03/2026
% =============================================================================

\chapter{Semana 04 (04/03/2026) — Estatística Descritiva II: Variabilidade e Dispersão}

% ---------------------------------------------------------------------------
\section{Objetivo da semana}
% ---------------------------------------------------------------------------

Olá pessoal! Na semana passada calculamos onde os dados estão (média, mediana,
moda). Hoje vamos responder a segunda pergunta fundamental: \textbf{quão
espalhados eles estão?} Um grupo com média R\$5.000 pode ter todo mundo
ganhando entre R\$4.800 e R\$5.200 — ou ter metade ganhando R\$1.000 e metade
ganhando R\$9.000. A média é a mesma; a \emph{história} é completamente diferente.

Dispersão é o que revela isso. E na primeira parte da aula, vocês vão apresentar
os \textbf{infográficos das funções Python} — a biblioteca da turma começa a crescer!

\begin{NoteBox}
\textbf{Atenção — aula compactada:} faltei na aula de quarta da semana passada,
então hoje vamos cobrir conteúdo acumulado. O ritmo será mais rápido que o
normal para garantir que tudo o que cai na Prova I esteja coberto.
\end{NoteBox}

\begin{BoardBox}
\textbf{Roteiro da aula (19h–22h)}
\begin{enumerate}
  \item 19h–19h40 — \textbf{Apresentações:} Infográficos de Função (ritmo acelerado)
  \item 19h40–20h10 — \textbf{CATCH-UP:} Covariância, Correlação de Pearson e Scatter Plot
  \item 20h10–20h40 — \textbf{TEORIA:} Amplitude, Variância e Desvio Padrão
  \item 20h40–21h10 — \textbf{TEORIA:} IQR, Regra dos 1,5×IQR e Boxplot
  \item 21h10–21h25 — \textbf{TEORIA:} Quando usar DP vs IQR + Comparação entre grupos
  \item 21h25–21h45 — \textbf{PRÁTICA EM DUPLA:} Tabelas 1 e 2 (10 linhas cada)
  \item 21h45–22h — \textbf{REVISÃO SEMANAL:} Questão tipo Prova I + lançamento do T2
\end{enumerate}
\end{BoardBox}

\begin{NoteBox}
\textbf{Leitura base:} Cap.\ 02 (Análise Univariada) — seções de variabilidade,
desvio padrão e boxplot. Livro: \textit{Estatística Prática para Cientistas de Dados} (Bruce).
\end{NoteBox}

% =============================================================================
\section{Parte I — Apresentações: Infográficos de Função (19h–20h)}
% =============================================================================

Pessoal, 2 minutos por aluno. Obrigatório mostrar no Jupyter ao vivo. Se o
código não rodar em sala, a análise crítica salva — mas o código precisa
estar funcional para a nota completa.

\begin{BoardBox}
\textbf{Checklist do professor durante as apresentações:}
\begin{tabular}{|p{4cm}|p{9cm}|}
\hline
\textbf{Ponto} & \textbf{O que observar} \\
\hline
Função correta      & Nome e assinatura exatamente conforme a lista \\
Docstring presente  & Parâmetros, retorno e exemplo descritos \\
Código comentado    & Ao menos uma linha de comentário por bloco lógico \\
Saída com dataset   & Screenshot ou live: rodando com dados reais da aula \\
Análise crítica     & Número + o que significa + limitação + recomendação \\
\hline
\end{tabular}
\end{BoardBox}

\begin{SolvedBox}
\textbf{Após as apresentações — consolidação coletiva:}

Vamos listar no quadro as funções apresentadas e criar o
\textbf{índice da biblioteca da turma}. A partir de hoje,
qualquer notebook da disciplina pode usar \texttt{from biblioteca\_turma import calcular\_media}
e evitar reescrever código toda semana. Isso é colaboração real de
equipe de dados.
\end{SolvedBox}

% =============================================================================
\section{Parte I-B — Catch-up: Covariância, Correlação e Scatter Plot (19h40–20h10)}
% =============================================================================

Pessoal, esse bloco é um adiantamento rápido da análise bivariada que será
aprofundada na Semana 06. Estou trazendo agora porque faltei na última quarta
e precisamos cobrir tudo que cai na Prova I. Aqui o objetivo é:
\textbf{entender o conceito}, \textbf{saber interpretar} e \textbf{distinguir
correlação de causalidade}.

% ---------------------------------------------------------------------------
\subsection{Scatter Plot — visão rápida}
% ---------------------------------------------------------------------------

\begin{BoardBox}
\textbf{O que é o scatter plot (para o quadro)}

O scatter plot coloca dois atributos nos eixos X e Y e plota cada observação
como um ponto. O padrão visual revela a relação entre as variáveis.

\textbf{O que observar:}
\begin{itemize}
  \item \textbf{Direção:} pontos sobem (positiva) ou descem (negativa)?
  \item \textbf{Intensidade:} concentrados ao longo de uma linha (forte) ou espalhados (fraca)?
  \item \textbf{Outliers:} pontos muito afastados do padrão geral.
  \item \textbf{Cor por categoria:} pintar por uma terceira variável pode revelar subgrupos.
\end{itemize}
\end{BoardBox}

% ---------------------------------------------------------------------------
\subsection{Covariância}
% ---------------------------------------------------------------------------

\begin{SolvedBox}
\textbf{O que é a covariância?} Mede o grau em que duas variáveis variam
conjuntamente. Se valores altos de $X$ acompanham valores altos de $Y$,
a covariância é positiva.

\begin{FormulaBox}{Covariância amostral}
\[
\mathrm{Cov}(X,Y) = \frac{1}{n-1}\sum_{i=1}^{n}(x_i - \bar{x})(y_i - \bar{y})
\]
\end{FormulaBox}

\textbf{Exemplo rápido (dados da Semana 03):}

\begin{center}
\begin{tabular}{|c|r|r|r|r|r|}
\hline
Vendedor & $x_i$ (Vendas) & $y_i$ (Meta) & $(x_i-\bar{x})$ & $(y_i-\bar{y})$ & Produto \\
\hline
Ana   & 4.200 & 4.000 & $-1.480$ & $-1.200$ & $1.776.000$ \\
Bruno & 6.000 & 5.500 & $+320$   & $+300$   & $96.000$    \\
Carla & 4.200 & 3.800 & $-1.480$ & $-1.400$ & $2.072.000$ \\
Diego & 9.800 & 8.000 & $+4.120$ & $+2.800$ & $11.536.000$ \\
Elisa & 4.200 & 4.700 & $-1.480$ & $-500$   & $740.000$   \\
\hline
\multicolumn{5}{|r|}{\textbf{Soma}} & $\mathbf{16.220.000}$ \\
\hline
\end{tabular}
\end{center}

$\bar{x} = 5.680$, $\bar{y} = 5.200$

\[
\mathrm{Cov}(X,Y) = \frac{16.220.000}{4} = 4.055.000
\]

Covariância positiva $\Rightarrow$ vendas e metas variam na mesma direção.

\textbf{Limitação:} o valor depende da escala. Usamos a \textbf{correlação} para normalizar.
\end{SolvedBox}

% ---------------------------------------------------------------------------
\subsection{Correlação de Pearson}
% ---------------------------------------------------------------------------

\begin{FormulaBox}{Correlação de Pearson ($r$)}
\[
r = \frac{\mathrm{Cov}(X,Y)}{s_X \cdot s_Y}, \quad r \in [-1, 1]
\]
\begin{itemize}
  \item $|r| < 0{,}3$: fraca \quad $0{,}3 \leq |r| < 0{,}7$: moderada \quad $|r| \geq 0{,}7$: forte
  \item $r > 0$: positiva (X sobe, Y sobe) \quad $r < 0$: negativa (X sobe, Y desce)
\end{itemize}
\end{FormulaBox}

\begin{BoardBox}
\textbf{REGRA DE OURO — Correlação $\neq$ Causalidade (para o quadro)}

Um $r$ alto mostra que as variáveis \textbf{andam juntas}, mas \textbf{não
prova} que uma causa a outra. Pode haver uma terceira variável (confundidor)
causando as duas.

\textbf{Exemplo clássico:} venda de sorvete e afogamentos têm $r > 0{,}7$.
Causa? Não — o \textbf{verão} causa os dois.

\textbf{Para provar causalidade:} precedência temporal + co-variação +
ausência de confundidor + (idealmente) experimento controlado.
\end{BoardBox}

% =============================================================================
\section{Parte II — Teoria: Medidas de Dispersão (20h10–21h25)}
% =============================================================================

% ---------------------------------------------------------------------------
\subsection{20h–20h30: Amplitude, Variância e Desvio Padrão}
% ---------------------------------------------------------------------------

Pessoal, medidas de posição nos dizem \emph{onde} os valores estão.
Medidas de dispersão nos dizem \emph{o quanto eles variam em torno desse centro}.
Vamos do mais simples para o mais robusto.

\subsubsection*{Amplitude}

\begin{BoardBox}
\textbf{Amplitude (para o quadro)}
\[
\text{Amplitude} = x_{\max} - x_{\min}
\]

\textbf{Exemplo:} salários = \{1.200, 4.500, 8.000, 15.000, 35.000\}
\[
\text{Amplitude} = 35.000 - 1.200 = R\$\;33.800
\]

\textbf{Vantagem:} cálculo imediato, fácil de comunicar.\\
\textbf{Problema grave:} depende exclusivamente dos dois extremos. Um único
outlier distorce completamente. Se o CEO ganhar R\$999.000, a amplitude sobe
para R\$997.800 — mas nada mudou para os outros quatro funcionários.\\
\textbf{Uso:} apenas como primeira verificação rápida. Nunca como única medida.
\end{BoardBox}

\subsubsection*{Variância}

\begin{SolvedBox}
\textbf{O que é a variância?} A variância mede a dispersão média dos desvios
quadráticos em relação à média. Elevamos ao quadrado para eliminar sinais
negativos e penalizar desvios grandes mais do que desvios pequenos.

\begin{FormulaBox}{Variância amostral}
\[
s^2 = \frac{\sum_{i=1}^{n}(x_i - \bar{x})^2}{n - 1}
\]
\end{FormulaBox}

\textbf{Por que $n-1$ e não $n$?} Porque estamos estimando a variância da
\emph{população} a partir de uma \emph{amostra}. Dividir por $n-1$ (Correção
de Bessel) corrige o viés que ocorre quando a amostra é pequena. Em
populações grandes isso faz pouca diferença; em amostras pequenas, faz muita.

\textbf{Problema da variância:} a unidade fica ao quadrado.
Variância de salários em R\$ sai em R\$² — impossível de interpretar diretamente.
É por isso que usamos o desvio padrão.
\end{SolvedBox}

\subsubsection*{Desvio Padrão}

\begin{BoardBox}
\textbf{Cálculo completo — Desvio Padrão de Vendas (quadro)}

Dados (Tabela 0 da semana passada): 4.200, 6.000, 4.200, 9.800, 4.200\\
Média: $\bar{x} = 5.680$

\begin{center}
\begin{tabular}{|r|r|r|}
\hline
$x_i$ & $x_i - \bar{x}$ & $(x_i - \bar{x})^2$ \\
\hline
4.200 & $-1.480$ & $2.190.400$ \\
6.000 & $+320$   & $102.400$   \\
4.200 & $-1.480$ & $2.190.400$ \\
9.800 & $+4.120$ & $16.974.400$ \\
4.200 & $-1.480$ & $2.190.400$ \\
\hline
\textbf{Soma} & & $23.648.000$ \\
\hline
\end{tabular}
\end{center}

\[
s^2 = \frac{23.648.000}{5-1} = 5.912.000
\qquad
s = \sqrt{5.912.000} \approx R\$\;2.432
\]

\textbf{Interpretação:} as vendas se afastam, em média, R\$2.432 da média
de R\$5.680. Diego (R\$9.800) e os três empatados em R\$4.200 criam essa
dispersão alta.
\end{BoardBox}

\begin{SolvedBox}
\textbf{Como interpretar o desvio padrão em contexto de negócio:}
\begin{itemize}
  \item \textbf{DP pequeno em relação à média:} dados homogêneos, previsíveis.
  Ex.: tempo de atendimento com média 9 min e DP 0,5 min → processo estável.
  \item \textbf{DP grande em relação à média:} dados heterogêneos, imprevisíveis.
  Ex.: tempo de atendimento com média 9 min e DP 12 min → há casos muito fora do padrão.
  \item \textbf{Comparação entre grupos com DP:} o grupo com menor DP é mais consistente,
  independentemente da média. Um time de vendas com média R\$5.000 e DP R\$500 é mais
  confiável que outro com média R\$6.000 e DP R\$3.000.
\end{itemize}
\end{SolvedBox}

Perguntem aos alunos: \textit{``Entre dois fornecedores com o mesmo preço médio,
qual vocês escolheriam: o com DP maior ou menor? Por quê isso muda dependendo
do contexto?''}

% ---------------------------------------------------------------------------
\subsection{20h30–21h: IQR, Regra dos 1,5×IQR e Boxplot}
% ---------------------------------------------------------------------------

\subsubsection*{Quartis e IQR}

\begin{BoardBox}
\textbf{Quartis e IQR (para o quadro)}

Quartis dividem os dados ordenados em quatro partes iguais:
\begin{itemize}
  \item $Q_1$ (1º quartil / P25): 25\% dos dados estão abaixo.
  \item $Q_2$ (mediana / P50): 50\% dos dados estão abaixo.
  \item $Q_3$ (3º quartil / P75): 75\% dos dados estão abaixo.
  \item $\text{IQR} = Q_3 - Q_1$: intervalo que contém os 50\% centrais dos dados.
\end{itemize}

\textbf{Cálculo com os dados de vendas:}
Ordenados: $4200,\; 4200,\; 4200,\; 6000,\; 9800$

\begin{itemize}
  \item $Q_1$: mediana da metade inferior $\{4200, 4200\} = 4200$
  \item $Q_2$: mediana geral $= 4200$
  \item $Q_3$: mediana da metade superior $\{6000, 9800\} = 7900$
  \item $\text{IQR} = 7900 - 4200 = R\$\;3.700$
\end{itemize}

\textbf{Quando $n$ é par,} a metade inferior e a superior são divididas
exatamente ao meio. Quando $n$ é ímpar, o valor central ($Q_2$) pode ser
incluído ou excluído dependendo da convenção — o Python usa interpolação.
\end{BoardBox}

\subsubsection*{Regra dos 1,5×IQR — Detecção formal de outliers}

\begin{SolvedBox}
\textbf{A regra de Tukey (1,5×IQR):}

\[
\text{Limite inferior} = Q_1 - 1{,}5 \times \text{IQR}
\qquad
\text{Limite superior} = Q_3 + 1{,}5 \times \text{IQR}
\]

Qualquer valor fora desses limites é classificado como \textbf{outlier}.

\textbf{Aplicando nos dados de vendas:}
\[
\text{LI} = 4200 - 1{,}5 \times 3700 = 4200 - 5550 = -1350
\]
\[
\text{LS} = 7900 + 1{,}5 \times 3700 = 7900 + 5550 = 13.450
\]

Nenhum valor está fora de $[-1350,\; 13450]$ $\Rightarrow$ sem outlier pela regra
de Tukey neste conjunto pequeno.

\textbf{Mas e se houvesse um erro de digitação de R\$98.000?}
$98.000 > 13.450$ $\Rightarrow$ \textbf{outlier confirmado} pela regra.

\textbf{Importante:} a regra sinaliza o valor para investigação — não ordena
a remoção automática.
\end{SolvedBox}

\subsubsection*{Boxplot — a visualização da dispersão}

\begin{BoardBox}
\textbf{Como ler um boxplot (para o quadro)}

\begin{center}
\begin{tabular}{|l|p{10cm}|}
\hline
\textbf{Elemento} & \textbf{O que representa} \\
\hline
Linha central da caixa & Mediana ($Q_2$) \\
Borda inferior da caixa & $Q_1$ (25\% dos dados) \\
Borda superior da caixa & $Q_3$ (75\% dos dados) \\
Altura da caixa & IQR (50\% centrais dos dados) \\
Bigodes (whiskers) & Dados dentro de $1{,}5 \times \text{IQR}$ \\
Pontos isolados & Outliers (além dos bigodes) \\
\hline
\end{tabular}
\end{center}

\textbf{O que um boxplot revela rapidamente:}
\begin{itemize}
  \item Caixa pequena + bigodes curtos $\Rightarrow$ dados concentrados, previsíveis.
  \item Caixa grande + bigodes longos $\Rightarrow$ dados muito dispersos.
  \item Mediana deslocada para cima da caixa $\Rightarrow$ assimetria negativa (cauda esquerda).
  \item Mediana deslocada para baixo da caixa $\Rightarrow$ assimetria positiva (cauda direita).
  \item Pontos isolados $\Rightarrow$ investigar outliers individualmente.
\end{itemize}
\end{BoardBox}

Perguntem aos alunos: \textit{``Quando o boxplot de um grupo A tem caixa menor
que o do grupo B mas mesma mediana, o que isso diz sobre a gestão de cada grupo?''}

% ---------------------------------------------------------------------------
\subsection{21h–21h20: Quando usar DP vs IQR + Comparação de grupos}
% ---------------------------------------------------------------------------

\begin{BoardBox}
\textbf{Guia de decisão: DP ou IQR? (para o quadro)}
\begin{tabular}{|l|l|l|}
\hline
\textbf{Situação} & \textbf{Use} & \textbf{Por quê} \\
\hline
Distribuição simétrica, sem outliers & Desvio Padrão & Usa todas as observações; mais eficiente \\
Há outliers ou assimetria clara & IQR & Robusto; não distorce com extremos \\
Comparação de grupos diferentes & Ambos & DP para magnitude; IQR para simetria \\
Comunicação executiva & IQR + Boxplot & Mais intuitivo para não-estatísticos \\
\hline
\end{tabular}
\end{BoardBox}

\begin{SolvedBox}
\textbf{Comparação de dispersão entre grupos — Exemplo:}

Dois times de suporte com mesmo tempo médio (9 min):
\begin{center}
\begin{tabular}{|l|c|c|c|c|}
\hline
\textbf{Time} & \textbf{Média} & \textbf{DP} & \textbf{IQR} & \textbf{Interpretação} \\
\hline
Time A & 9 min & 0,8 min & 1,2 min & Muito consistente — processo padronizado \\
Time B & 9 min & 6,5 min & 8,0 min & Muito variável — há casos extremamente rápidos e lentos \\
\hline
\end{tabular}
\end{center}

\textbf{Decisão de negócio:} se a meta é previsibilidade (SLA), Time A é melhor.
Se o objetivo é ter \emph{alguns} atendentes muito rápidos (para triagem), Time B pode
ser valioso — mas precisa de investigação dos casos lentos.
\end{SolvedBox}

% =============================================================================
\section{Parte III — Prática em Dupla (21h20–21h45)}
% =============================================================================

\begin{SolvedBox}
\textbf{Instruções:}
\begin{enumerate}
  \item Calcule Amplitude, Variância, Desvio Padrão para a coluna numérica.
  \item Calcule Q1, Q3 e IQR.
  \item Aplique a regra de Tukey — há outlier?
  \item Compare two grupos usando DP e IQR.
  \item Escreva 1 frase de interpretação para cada resultado.
\end{enumerate}
\end{SolvedBox}

\begin{BoardBox}
\textbf{Tabela 1 — Tempo de Entrega por Transportadora (calcule dispersão)}

\medskip
\begin{tabular}{|c|l|r|l|}
\hline
\textbf{ID} & \textbf{Pedido} & \textbf{Dias para entrega} & \textbf{Transportadora} \\
\hline
01 & Pedido A & 3  & LogRápida \\
02 & Pedido B & 5  & FreteJá  \\
03 & Pedido C & 4  & LogRápida \\
04 & Pedido D & 3  & LogRápida \\
05 & Pedido E & 21 & FreteJá  \\
06 & Pedido F & 4  & FreteJá  \\
07 & Pedido G & 3  & LogRápida \\
08 & Pedido H & 6  & FreteJá  \\
09 & Pedido I & 4  & LogRápida \\
10 & Pedido J & 5  & FreteJá  \\
\hline
\end{tabular}

\smallskip
\textbf{Calcule (coluna Dias para entrega):} Amplitude · Média · Mediana · Variância · DP · Q1 · Q3 · IQR · Limites Tukey

\textbf{Compare:} média e DP de cada Transportadora separadamente.

\smallskip
\textbf{Respostas esperadas:}
\begin{itemize}
  \item Soma: $3+5+4+3+21+4+3+6+4+5 = 58$ \quad \textbf{Média:} $58 \div 10 = 5{,}8$ dias
  \item Ordenados: $3,3,3,4,4,4,5,5,6,21$ \quad \textbf{Mediana:} $\frac{4+4}{2} = 4$ dias
  \item $Q_1 = 3$; $Q_3 = 5$; $\text{IQR} = 2$ \quad \textbf{Amplitude:} $21-3=18$
  \item $\text{LI} = 3 - 3 = 0$; $\text{LS} = 5 + 3 = 8$ $\Rightarrow$ \textbf{21 dias é outlier!}
  \item Variância: $\approx 27{,}5$; DP $\approx 5{,}2$ dias (distorcido pelo outlier)
  \item \textbf{LogRápida:} média $= (3+4+3+3+4)/5 = 3{,}4$ dias, DP $\approx 0{,}5$
  \item \textbf{FreteJá:} média $= (5+21+4+6+5)/5 = 8{,}2$ dias, DP $\approx 6{,}7$
  \item \textbf{Decisão:} FreteJá tem outlier real (21 dias — provavelmente extravio);
  LogRápida é 2,4× mais rápida E 13× mais consistente (DP). Contratar LogRápida.
\end{itemize}
\end{BoardBox}

\begin{BoardBox}
\textbf{Tabela 2 — Pontuação de Clientes por Avaliação (calcule dispersão)}

\medskip
\begin{tabular}{|c|l|r|l|}
\hline
\textbf{ID} & \textbf{Cliente} & \textbf{NPS (0–100)} & \textbf{Plano} \\
\hline
01 & Cliente 01 & 85 & Premium \\
02 & Cliente 02 & 42 & Básico  \\
03 & Cliente 03 & 90 & Premium \\
04 & Cliente 04 & 78 & Premium \\
05 & Cliente 05 & 15 & Básico  \\
06 & Cliente 06 & 88 & Premium \\
07 & Cliente 07 & 60 & Básico  \\
08 & Cliente 08 & 91 & Premium \\
09 & Cliente 09 & 35 & Básico  \\
10 & Cliente 10 & 83 & Premium \\
\hline
\end{tabular}

\smallskip
\textbf{Calcule:} Média · Mediana · DP · Q1 · Q3 · IQR · Limites Tukey para NPS total.\\
\textbf{Compare:} média e DP de NPS por Plano (Premium vs Básico).

\smallskip
\textbf{Respostas esperadas:}
\begin{itemize}
  \item Soma: $85+42+90+78+15+88+60+91+35+83 = 667$ \quad \textbf{Média:} $66{,}7$
  \item Ordenados: $15,35,42,60,78,83,85,88,90,91$ \quad \textbf{Mediana:} $(78+83)/2 = 80{,}5$
  \item $Q_1 = 42$; $Q_3 = 88$; $\text{IQR} = 46$; $\text{LI} = 42 - 69 = -27$; $\text{LS} = 88 + 69 = 157$
  \item Sem outliers; DP $\approx 27{,}4$ — alta dispersão para uma métrica de 0–100
  \item \textbf{Premium:} média $= (85+90+78+88+91+83)/6 = 85{,}8$; DP $\approx 4{,}7$
  \item \textbf{Básico:} média $= (42+15+60+35)/4 = 38{,}0$; DP $\approx 18{,}8$
  \item \textbf{Insight:} diferença de 47,8 pontos de NPS entre planos; Básico tem 4×
  mais variabilidade — experiência inconsistente. Média geral (66,7) mascara os dois grupos.
\end{itemize}
\end{BoardBox}

% =============================================================================
\section{Parte IV — Atividade Individual Final (21h45–22h)}
% =============================================================================

Pessoal, última parte da noite — individual, sem consulta ao colega.
Apliquem tudo de hoje: amplitude, DP, IQR, regra de Tukey e comparação entre grupos.

\begin{SolvedBox}
\textbf{Instruções:}
\begin{enumerate}
  \item Calcule Amplitude, Variância, Desvio Padrão para a coluna Faturamento.
  \item Calcule Q1, Q3 e IQR. Aplique a regra de Tukey — há outlier?
  \item Compare Faturamento por Região: calcule média E DP de cada grupo.
  \item Identifique qual região é mais consistente e qual tem maior dispersão.
  \item Redija 1 insight: \emph{``O dado mostra X, que significa Y, portanto Z.''
}\end{enumerate}
\textbf{Tempo:} 15 minutos. Entrega: folha ou foto no AVA.
\end{SolvedBox}

\begin{BoardBox}
\textbf{Tabela 3 — Faturamento Mensal por Representante}

\medskip
\begin{tabular}{|c|l|r|l|}
\hline
\textbf{ID} & \textbf{Representante} & \textbf{Faturamento (R\$ mil)} & \textbf{Região} \\
\hline
01 & Rep 01 & 42 & Norte \\
02 & Rep 02 & 38 & Sul   \\
03 & Rep 03 & 45 & Norte \\
04 & Rep 04 & 41 & Norte \\
05 & Rep 05 & 120 & Sul  \\
06 & Rep 06 & 39 & Sul   \\
07 & Rep 07 & 43 & Norte \\
08 & Rep 08 & 37 & Sul   \\
09 & Rep 09 & 44 & Norte \\
10 & Rep 10 & 40 & Sul   \\
\hline
\end{tabular}
\end{BoardBox}

\begin{SolvedBox}
\textbf{Gabarito (professor — não mostrar antes da correção):}

\textbf{Geral:}
\begin{itemize}
  \item Soma: $42+38+45+41+120+39+43+37+44+40 = 489$ mil
  \item \textbf{Média:} $48{,}9$ mil \quad Ordenados: $37,38,39,40,41,42,43,44,45,120$
  \item \textbf{Mediana:} $(41+42)/2 = 41{,}5$ mil \quad $Q_1 = 39$; $Q_3 = 44$; $\text{IQR} = 5$
  \item $\text{LS} = 44 + 7{,}5 = 51{,}5$ $\Rightarrow$ \textbf{Rep 05 (R\$120 mil) é outlier}
  \item DP $\approx 24{,}8$ — fortemente inflado pelo outlier
\end{itemize}

\textbf{Por região:}
\begin{itemize}
  \item \textbf{Norte} (Rep 01,03,04,07,09): média $= (42+45+41+43+44)/5 = 43$; DP $\approx 1{,}6$
  \item \textbf{Sul} (Rep 02,05,06,08,10): média $= (38+120+39+37+40)/5 = 54{,}8$; DP $\approx 31{,}9$
\end{itemize}

\textbf{Insight esperado:} ``O dado mostra que o Sul tem faturamento médio
R\$11,8 mil acima do Norte, mas um DP 20× maior (31,9 vs 1,6). Isso significa
que o resultado do Sul depende de um único representante (Rep 05, R\$120 mil).
Sem ele, o Sul médio seria R\$38,5 mil — abaixo do Norte. Portanto, recomendo
mapear a carteira do Rep 05 como risco de concentração e capacitar os demais.''
\end{SolvedBox}

% =============================================================================
\section{Revisão Semanal — Questão Tipo Prova I}
% =============================================================================

\begin{SolvedBox}
\textbf{Questão de Revisão (Semana 04) — Mediana, Outliers e Dispersão}

Uma rede de supermercados com 8 lojas na Região Metropolitana do Recife coletou
o faturamento mensal (em R\$) do setor de hortifrúti em janeiro de 2026:

\begin{center}
\begin{tabular}{|l|l|r|}
\hline
\textbf{Loja} & \textbf{Localização} & \textbf{Faturamento (R\$)} \\
\hline
L1 & Boa Viagem  & 45.000 \\
L2 & Casa Forte  & 52.000 \\
L3 & Olinda      & 38.000 \\
L4 & Paulista    & 41.000 \\
L5 & Piedade     & 39.000 \\
L6 & Jaboatão    & 36.000 \\
L7 & Camaragibe  & 35.000 \\
L8 & Derby       & 142.000 \\
\hline
\end{tabular}
\end{center}

O estagiário calculou a média aritmética: R\$53.500. O analista sênior argumentou
que a mediana é mais representativa. Além disso:
$Q_1 = 37.000$, $Q_3 = 48.500$, $\text{IQR} = 11.500$.

\textbf{Pergunta:} Calcule a mediana, verifique se L8 é outlier pelo critério
$1{,}5 \times \text{IQR}$, e justifique a escolha entre média e mediana.

\medskip
\textbf{Resolução:}

Ordenados: $35.000,\; 36.000,\; 38.000,\; 39.000,\; 41.000,\; 45.000,\; 52.000,\; 142.000$

$n = 8$ (par) $\Rightarrow$ Mediana $= \dfrac{39.000 + 41.000}{2} = \mathbf{R\$\;40.000}$

Limite superior: $Q_3 + 1{,}5 \times \text{IQR} = 48.500 + 17.250 = \mathbf{R\$\;65.750}$

$142.000 > 65.750$ $\Rightarrow$ \textbf{L8 é outlier.}

\textbf{Justificativa:} A mediana (R\$40.000) é preferível porque é resistente a
valores extremos e representa melhor o faturamento típico das 7 lojas ``normais''.
A média (R\$53.500) é inflada pelo outlier L8 e não representa nenhuma loja
individualmente.
\end{SolvedBox}

% ---------------------------------------------------------------------------
\section{Entrega — Trabalho 2 (lance em sala)}
% ---------------------------------------------------------------------------

\begin{SolvedBox}
\textbf{Trabalho 2 — Mini-relatório AED (entrega: 11/03/2026 às 23h59 via AVA):}
\begin{itemize}
  \item[$\square$] Dataset próprio ou do Kaggle com pelo menos 50 linhas e 3 colunas numéricas
  \item[$\square$] Calcule Média, Mediana, Moda, DP e IQR para cada coluna numérica
  \item[$\square$] Identifique e investigue outliers pela regra de Tukey
  \item[$\square$] Compare pelo menos 2 grupos usando média + DP
  \item[$\square$] 2 visuais: 1 boxplot + 1 histograma (Python, Tableau ou Excel)
  \item[$\square$] Big Idea (1 frase): o que o dado revela e que decisão ele apoia
\end{itemize}
\end{SolvedBox}

% ---------------------------------------------------------------------------
\section{Materiais Preparados}
% ---------------------------------------------------------------------------
\begin{itemize}
  \item Slides: fórmulas animadas passo a passo (variância → DP → IQR → boxplot)
  \item Tabelas 1, 2 e 3 impressas (uma por dupla/individual)
  \item Dataset da aula disponível no AVA (\texttt{semana04\_vendas.csv})
  \item Notebook base com tabelas carregadas + esqueleto das funções do aluno
\end{itemize}

% ---------------------------------------------------------------------------
\section{Próximo passo}
% ---------------------------------------------------------------------------

Na \textbf{Semana 05} vamos olhar para o \textbf{formato} dos dados —
histogramas, simetria, assimetria e caudas. Também vamos enfrentar o problema
mais chato da vida de um analista: os \textbf{dados ausentes} (missing).
Como identificar, como decidir o que fazer e como documentar a decisão.

Perguntem aos alunos: \textit{``Se vocês fossem calcular o DP dos salários
de uma empresa com 1 CEO e 99 funcionários de base, qual seria o resultado?
O DP seria um número útil para apresentar ao RH? O que vocês usariam no lugar?''}
